\section{Learning from first principles}\label{sec:fp}

This section derives the training method step by step. Each step ends with a pointer to the section that
states the exact form used in the code. Nothing here is specific to Tigers \& Goats except the last two
subsections.

\subsection{The game as a decision process}\label{sec:fp-mdp}

Training controls one side, the \emph{learner}. The other side is played by a fixed opponent policy $\mu$
(Section~\ref{old:2.4}). From the learner's point of view the opponent is part of the world: the learner
picks an action $a_t$, the world applies it, lets the opponent reply, and returns the next observation and
a reward. That is a \textbf{Markov decision process}
\[
\mathcal{M}_\mu = \big(\mathcal{S},\ \mathcal{A},\ P_\mu,\ r,\ \gamma\big),
\qquad s_{t+1} \sim P_\mu(\cdot \mid s_t, a_t), \qquad r_t = r(s_t, a_t, s_{t+1}),
\]
with states $\mathcal{S}$, the 115 actions $\mathcal{A}$, transition law $P_\mu$ (which contains the opponent and
its randomness), reward $r$ and discount $\gamma \in (0, 1]$. The subscript matters: a different opponent is a
different decision process.

Strictly, the observation $s_t$ of Section~\ref{old:2.1} is not the full state. It leaves out the turn
counter and the repeat history, which affect rewards and endings, so the process is partially observed. The
policies here ignore that and act on $s_t$ alone.

A \textbf{policy} $\pi_\theta(a \mid s)$ is a probability distribution over actions, computed by a neural network
with parameters $\theta$ (Section~\ref{old:3}). Playing $\pi_\theta$ in $\mathcal{M}_\mu$ produces a random
\emph{trajectory} $\xi = (s_0, a_0, r_0, s_1, a_1, r_1, \dots, s_T)$ with probability
\begin{equation}\label{eq:traj}
p_\theta(\xi) = p_0(s_0) \prod_{t=0}^{T-1} \pi_\theta(a_t \mid s_t)\, P_\mu(s_{t+1} \mid s_t, a_t).
\end{equation}
The \textbf{return} from step $t$ and the \textbf{objective} are
\begin{equation}\label{eq:objective}
G_t = \sum_{l \ge 0} \gamma^{l}\, r_{t+l},
\qquad
J(\theta) = \E_{\xi \sim p_\theta}\big[G_0\big].
\end{equation}
Training means finding $\theta$ that makes $J(\theta)$ large. Three functions describe how good a policy $\pi$
is from a given point:
\[
V^\pi(s) = \E_\pi[G_t \mid s_t = s], \qquad
Q^\pi(s, a) = \E_\pi[G_t \mid s_t = s, a_t = a], \qquad
A^\pi(s, a) = Q^\pi(s, a) - V^\pi(s).
\]
$V^\pi$ is the expected return from $s$; $Q^\pi$ is the same after first taking $a$; the \textbf{advantage} $A^\pi$
says how much better $a$ is than the policy's average choice. They satisfy the Bellman equation
$V^\pi(s) = \E_{a \sim \pi,\, s' \sim P_\mu}\big[r + \gamma V^\pi(s')\big]$, with $V^\pi = 0$ after a terminal state.

\subsection{The policy gradient}\label{sec:fp-pg}

$J$ is an expectation over trajectories whose distribution depends on $\theta$. Its gradient follows from the
identity $\nabla_\theta p_\theta = p_\theta \nabla_\theta \log p_\theta$:
\[
\nabla_\theta J(\theta) = \int \nabla_\theta p_\theta(\xi)\, G_0(\xi)\, d\xi = \E_{\xi \sim p_\theta}\big[\nabla_\theta \log p_\theta(\xi)\ G_0(\xi)\big].
\]
In $\log p_\theta(\xi)$, from \eqref{eq:traj}, only the policy factors depend on $\theta$. The start
distribution, the game rules and the opponent drop out:
\begin{equation}\label{eq:pg-raw}
\nabla_\theta J(\theta) = \E\Big[\sum_{t} \nabla_\theta \log \pi_\theta(a_t \mid s_t)\ G_0\Big].
\end{equation}
So the gradient can be estimated from played games alone, without a model of the rules or the opponent.
Two refinements reduce the noise of \eqref{eq:pg-raw} without changing its expectation.

\textbf{Causality.} An action cannot affect rewards that came before it. The rewards $r_0, \dots, r_{t-1}$
multiply $\nabla \log \pi_\theta(a_t \mid s_t)$ with expectation zero, so $G_0$ can be replaced by
$\gamma^t G_t$.

\textbf{Baseline.} For any function $b(s)$ that does not depend on the action,
\[
\E_{a \sim \pi_\theta(\cdot \mid s)}\big[\nabla_\theta \log \pi_\theta(a \mid s)\ b(s)\big]
 = b(s)\, \nabla_\theta \sum_{a} \pi_\theta(a \mid s) = b(s)\, \nabla_\theta 1 = 0 .
\]
Subtracting $b(s_t)$ from the return therefore leaves the gradient unbiased. The choice $b = V^\pi$ gives the
\textbf{policy gradient theorem} in the form used here:
\begin{equation}\label{eq:pg}
\nabla_\theta J(\theta) = \E\Big[\sum_{t} \gamma^t\, \nabla_\theta \log \pi_\theta(a_t \mid s_t)\ A^{\pi_\theta}(s_t, a_t)\Big].
\end{equation}
Reading: raise the probability of actions that did better than expected, lower it for those that did
worse, in proportion to how much. Like most implementations, the code drops the factor $\gamma^t$ and
weights all time steps equally.

\subsection{Estimating the advantage}\label{sec:fp-gae}

$A^\pi$ is unknown. A second network, the \textbf{critic} $V_\psi(s)$, learns to predict the return, and the
advantage is estimated from it. The \textbf{temporal-difference residual}
\[
\delta_t = r_t + \gamma\, V_\psi(s_{t+1}) - V_\psi(s_t)
\]
compares the critic's prediction before the step with the reward actually received plus the prediction
after it. If the critic were exact, $\E[\delta_t \mid s_t, a_t] = A^\pi(s_t, a_t)$. With an imperfect critic
$\delta_t$ is biased, but its variance is small because it uses one random step. Looking $j$ steps ahead
before trusting the critic gives the estimators
\[
\hat A^{(j)}_t = \sum_{l=0}^{j-1} \gamma^l\, \delta_{t+l} = r_t + \gamma r_{t+1} + \dots + \gamma^{j-1} r_{t+j-1} + \gamma^{j} V_\psi(s_{t+j}) - V_\psi(s_t).
\]
Small $j$ leans on the critic (bias); large $j$ leans on the sampled rewards (variance).
\textbf{Generalized advantage estimation} averages them with weights $(1-\lambda)\lambda^{j-1}$:
\begin{equation}\label{eq:gae}
\hat A_t = (1 - \lambda) \sum_{j \ge 1} \lambda^{j-1} \hat A^{(j)}_t = \sum_{l \ge 0} (\gamma\lambda)^l\, \delta_{t+l},
\qquad\text{computed backwards as}\quad \hat A_t = \delta_t + \gamma\lambda\, \hat A_{t+1}.
\end{equation}
$\lambda = 0$ is the one-step estimate and $\lambda = 1$ the full sampled return minus the baseline. All
models here use $\gamma = 0.99$ and $\lambda = 0.95$. The critic's regression target is
$\hat R_t = \hat A_t + V_\psi(s_t)$.

\textbf{Termination and truncation.} At a real ending (a win or a loss) nothing follows, so the value after it
is zero and the recursion \eqref{eq:gae} stops. At a \emph{truncation} (the turn limit) the game was merely cut
off, so the critic's value of the last state stands in for the missing future. Which endings count as which
has consequences. If an illegal move is reported as a truncation, the bootstrap adds back almost exactly the
value the state had before the move, and the penalty for the illegal move nearly cancels. Section~\ref{old:2.3}
and Section~\ref{old:5.4} return to this. The exact formulas with episode boundaries are in Section~\ref{old:4}.

\subsection{From a gradient to a safe update: PPO}\label{sec:fp-ppo}

The gradient \eqref{eq:pg} is valid only at the parameters that generated the data. One small step per batch
of games would waste most of the data; many steps on the same batch can move the policy so far that the
data no longer describes it. Proximal policy optimization takes several steps while keeping the new policy
close to the old one.

\textbf{Reusing data.} Let $\theta_{\text{old}}$ be the parameters that played the games. For a candidate
$\theta$, importance sampling re-weights each recorded action by the \textbf{probability ratio}
\[
\rho_t(\theta) = \frac{\pi_\theta(a_t \mid s_t)}{\pi_{\theta_{\text{old}}}(a_t \mid s_t)},
\qquad
S(\theta) = \E_{\text{old}}\big[\rho_t(\theta)\, \hat A_t\big].
\]
At $\theta = \theta_{\text{old}}$ the ratio is 1 and $\nabla_\theta S = \E[\nabla_\theta \log \pi_\theta\, \hat A_t]$, the policy
gradient. So maximizing the \emph{surrogate} $S$ starts in the right direction.

\textbf{Why it must stay close.} The exact change in the objective is (Kakade and Langford)
\[
J(\theta) - J(\theta_{\text{old}}) = \E_{\xi \sim \pi_\theta}\Big[\sum_t \gamma^t A^{\pi_{\text{old}}}(s_t, a_t)\Big],
\]
an expectation over the states the \emph{new} policy visits. $S(\theta)$ uses the states the old policy
visited. The two agree to first order, and the gap grows with the distance between the policies. Trust-region
methods bound that distance by the KL divergence. PPO enforces it more simply.

\textbf{The clipped objective.} With clip range $\epsilon$ (0.2 here), PPO maximizes
\begin{equation}\label{eq:clip}
S^{\text{clip}}(\theta) = \E_{\text{old}}\Big[\min\big(\rho_t \hat A_t,\ \operatorname{clip}(\rho_t,\, 1 - \epsilon,\, 1 + \epsilon)\, \hat A_t\big)\Big].
\end{equation}
Take the two cases. If $\hat A_t > 0$ the action was good and the term is $\min(\rho_t, 1 + \epsilon)\, \hat A_t$:
raising its probability is rewarded only up to a factor $1 + \epsilon$. If $\hat A_t < 0$ the term is
$\max(\rho_t, 1 - \epsilon)\, \hat A_t$: lowering its probability is rewarded only down to a factor $1 - \epsilon$.
Beyond those points the gradient is zero, so there is no incentive to move further. Moves in the harmful
direction are never clipped. $S^{\text{clip}}$ is thus a pessimistic lower bound on $S$ that removes the
reward for leaving the trust region. Section~\ref{old:5.3} lists two alternatives used in the experiments, a
KL penalty and the smooth SPO penalty.

\textbf{The complete loss.} Three terms are minimized together on each minibatch $\mathcal{B}$:
\begin{equation}\label{eq:total}
\mathcal{J}(\theta, \psi) =
\underbrace{-\,S^{\text{clip}}(\theta)}_{\mathcal{J}^\pi:\ \text{improve the policy}}
\ +\ c_{\text{vf}} \underbrace{\E_{\mathcal{B}}\big[(\hat R_t - V_\psi(s_t))^2\big]}_{\mathcal{J}^V:\ \text{fit the critic}}
\ +\ c_{\text{ent}} \underbrace{\big(-\E_{\mathcal{B}}\,\mathcal{H}[\pi_\theta(\cdot \mid s_t)]\big)}_{\mathcal{J}^{\mathcal{H}}:\ \text{keep exploring}} .
\end{equation}
The entropy $\mathcal{H}[\pi] = -\sum_a \pi(a) \log \pi(a)$ is largest for a uniform policy. Rewarding it stops
the policy from committing to one move before it has tried the others. Algorithm~\ref{alg:ppo} puts the
pieces together.

\begin{algorithm}[H]
\caption{One training run of PPO (the form used for every model; Sections~\ref{old:4}--\ref{old:5})}\label{alg:ppo}
\begin{tabbing}
\hspace{1.6em}\=\hspace{1.6em}\=\hspace{1.6em}\=\kill
\textbf{repeat} until the step budget is used: \\
\> \textbf{collect:} play $N_{\text{steps}}$ steps in each of $N_{\text{env}}$ games with $\pi_{\theta_{\text{old}}}$; store $s_t, a_t, r_t$, $V_\psi(s_t)$, $\log \pi_{\theta_{\text{old}}}(a_t \mid s_t)$ \\
\> \textbf{estimate:} compute $\hat A_t$ backwards by \eqref{eq:gae} and $\hat R_t = \hat A_t + V_\psi(s_t)$ \\
\> \textbf{for} epoch $= 1, \dots, N_{\text{ep}}$: \\
\>\> shuffle the $N_{\text{env}} N_{\text{steps}}$ transitions and split them into minibatches of size $B$ \\
\>\> \textbf{for} each minibatch $\mathcal{B}$: \\
\>\>\> normalize the advantages in $\mathcal{B}$ to mean 0 and standard deviation 1 \\
\>\>\> take one Adam step on $\mathcal{J}(\theta, \psi)$ of \eqref{eq:total}, with the gradient norm clipped to $g_{\max}$ \\
\> $\theta_{\text{old}} \leftarrow \theta$
\end{tabbing}
\end{algorithm}

\subsection{Illegal moves: masking and what it does not teach}\label{sec:fp-mask}

Most of the 115 actions are illegal in any given position. Two ways to deal with that lead to very
different networks.

\textbf{Masking} removes illegal actions before sampling: the masked policy is the raw policy restricted to the
legal set $\mathcal{L}(s)$ and renormalized, $\pi^m_\theta(a \mid s) = \pi_\theta(a \mid s)\, \ind[a \in \mathcal{L}(s)] / P_{\mathcal{L}}(s)$.
All of the above applies with $\pi^m_\theta$ in place of $\pi_\theta$. Its log-probability has the gradient
\[
\frac{\partial \log \pi^m_\theta(a \mid s)}{\partial z_b} =
\begin{cases} \ind[a = b] - \pi^m_\theta(b \mid s) & b \in \mathcal{L}(s) \\ 0 & b \notin \mathcal{L}(s) \end{cases}
\]
with respect to the logits $z$. The logit of an illegal action receives \emph{no gradient from any term of
the loss}. The network is never told that those actions are bad; the mask does that job at play time, and
the network is useless without it (Section~\ref{old:9.3}, finding 1).

\textbf{No mask} lets the policy sample any action and makes the environment punish illegal ones. The reward
signal is weak for this purpose: it arrives once per mistake and competes with everything else the reward
encodes.

\textbf{The legality loss} adds direct supervision. The mask is known at every visited state, so it can serve
as a label for the event ``the chosen action is legal'':
\[
\mathcal{J}^{\text{legal}}_{\text{nll}} = -\E_{\mathcal{B}}\big[\log P_{\mathcal{L}}(s_t)\big],
\qquad P_{\mathcal{L}}(s) = \sum_{a \in \mathcal{L}(s)} \pi_\theta(a \mid s).
\]
It is added to \eqref{eq:total} with weight $c_{\text{legal}}$. Unlike the reward, it gives a gradient at every
state and on every logit: for an illegal action $b$, $\partial \mathcal{J}^{\text{legal}}_{\text{nll}} / \partial z_b = \pi_\theta(b \mid s) > 0$,
which pushes that logit down. Section~\ref{old:5.4} gives the details and a variant that fails, and
Section~\ref{old:9} measures the effect.

\subsection{What a trained model is, and is not}\label{sec:fp-br}

The objective \eqref{eq:objective} is defined inside $\mathcal{M}_\mu$. Its maximizer is a \textbf{best response}
to the opponent $\mu$ under the shaped rewards of Appendix~\ref{app:rewards}. It is not the game-theoretic
optimum, which must hold against \emph{every} opponent. Two consequences follow, and Part~IV measures both.

\begin{itemize}
\item A best response may win more often against $\mu$ than a perfect player does, because it steers into
  the positions where $\mu$ goes wrong.
\item A best response may lose to an opponent it never met, in positions it never visited.
\end{itemize}

Part~III computes the other object: the value of every position against the best possible opponent.
